


A $t$-spread of $\PG(n,q)$ is a set of disjoint 
$\PG(t,q)$ that cover all of $\PG(n,q)$ pointwise. 
$t$-spreads in $\PG(n,q)$ exist if and only if $t+1$ divides $n+1.$


Table~\ref{tab:spreads:define} lists the Orbiter commands to create spreads.
\orbitertableofcommands{spread_create}


Table~\ref{tab:spreads:activity} lists the Orbiter spread activities.
\orbitertableofcommands{spread_activity}



Table~\ref{tab:spread:table} lists the Orbiter commands to create a table of spreads.
\orbitertableofcommands{spread_table}

Table~\ref{tab:spread:table:activity} lists the Orbiter activities for spread tables.
\orbitertableofcommands{spread_table_activity}




Table~\ref{tab:spreads:classify} list the Orbiter commands to create an object that can classify spreads.
\orbitertableofcommands{spread_classify}


Table~\ref{tab:spreads:classify:activity} 
lists the Orbiter commands that apply to an object that can classify spreads.
\orbitertableofcommands{spread_classify_activity}




\bigskip

The following two commands create the two spreads of order 9, 
relying on the Orbiter knowledge base.

\sourcecode{create_spread_9a}


\sourcecode{create_spread_9b}

The first spread is the Desarguesian spread, 
with automorphism group of order 5760.
The second spread is the Hall spread 
with automorphism group of order 1920.

\bigskip

Spreads can be defined using spread sets. 
A spread set is a set of $q^k$ matrices 
of size $k \times k$ over $\bbF_q$ such that $A_i-A_j$ 
is nonsingular for all $i \neq j.$ 
Let us look at an example. 
The spread due to Rao and Rao~\cite{RaoRao1976} 
can be defined using the following makefile variable. 

\sourcecodevariable{SPREAD_SET_27_RAO_RAO}

Each line represents one matrix of the spread set, 
with matrix entries being listed consecutively.
The following command can be used to define the spread:

\sourcecode{create_spread_Rao_Rao_27}




\bigskip


The following command creates the Desarguesian line-spread in $\PG(3,2)$: 
% (so $s=2,$ $t=s-1=1$, $m=2$, $q=2$, and $Q=4$):

\sourcecode{desarguesian_spread_in_PG_3_2}

The cheat sheet contains the following spread:\\
\orbiteroutput{GRAPHICS/LATEX/desarguesian_spread_PG_3_2}

\bigskip

The following command creates the Desarguesian plane-spread in $\PG(5,2)$:

\sourcecode{desarguesian_spread_in_PG_5_2}

\orbiteroutput{GRAPHICS/LATEX/desarguesian_spread_PG_5_2}

\bigskip


A spread table can be used to enumerate labeled spreads. 
For instance, we can create all labeled spreads associated with one particular isomorphism type. 
This means we can create spreads by orbits. We can choose the isomorphism type of spreads for which we want the orbits computed. 
In the following example, we create the labeled spreads associated with the unique spread in $\PG(3,2).$ 
There are exactly 56 spreads in the orbit:

\sourcecode{PG_3_2_spread_table}

The spreads are stored in a file 
$$
\verb'SPREAD_TABLES_2/spread_4_table.csv'
$$
Here, the integer 4 in the file name represents the order of the translation plane associated with the spreads. 


%Apart from the second spread element, 
%the left halves of the generator matrices of the subspaces in the Desarguesian spread 
%are the elements of $\bbF_8$ in a matrix representation over $\bbF_2$.

\bigskip

Two $t$-spreads are isomorphic if there is a collineation which maps one to the other. 
The classification problem for $t$-spreads 
is the problem of determining a complete set of pairwise non-isomorphic 
$t$-spreads. 
Orbiter can be used to classify spreads for small parameters. 
To enhance the algorithm, 
the method of classification by substructure is used.
Classification by substructure is a method to classify objects.
In this method, we do not classify the whole poset. 
Instead, we classify the substructures, and then move directly to the 
desired objects, thereby skipping over intermediate objects.
The jump over intermediate objects is facilitated 
by combinatorial searching, 
which may process nodes faster that classification with isomorph rejection. 
The combinatorial search is facilitated by an auxiliary algorithm. 
This auxiliary algorithm can be clique finding, exact cover, or any other computer science primitives. 
Table~\ref{tab:isomorph} shows the commands available for classification by substructure.
\orbitertableofcommands{isomorph}

\bigskip

Let us look at some examples. 



\bigskip

At first, we look at an example which is sufficiently small and can be 
solved using the standard method. 
Here, the standard method is poset classification algorithm for partial spreads.
Suppose we want to classify the line spreads in $\PG(3,4)$ under the action of $\PGGL(4,4).$ 
Under the Andr\'e, Bruck-Bose construction~\cite{Andre54,BruckBose64}, 
these spreads correspond to translation planes of order $16$ with kernel $\bbF_4.$
In order to classify the spreads of $\PG(3,4)$, we utilize poset classification, 
as shown in the following command

\sourcecode{classify_spreads_16_4}

The algorithm computes the poset of orbits for the group $G =\PGGL(4,4)$ acting on 
the poset of partial spreads in $\PG(3,4)$. 
The poset of orbits is shown below:

\orbiteroutput{GRAPHICS/WRAPPERS/spreads_16_4_draw}

Up to isomorphism, there are exactly 
three line-spreads in $\PG(3,4)$ (corresponding 
to the three nodes at the bottom of the poset of orbits in the figure). 
These three spreads are the Desarguesian spread, 
the Hall spread, and the semifield spread, respectively. 
Here is the relevant output taken from the latex report:
\orbiteroutput{GRAPHICS/LATEX/partial_spreads_PG_3_4}
The three spreads in $\PG(3,4)$ can be distinguished by their stabilizer orders.


%Table~\ref{tab:spreads:3:4} lists the line spreads in $\PG(3,4)$ 
%according to their orbiter catalogue number (OCN).

%\input GRAPHICS/TABLES/spreads_in_PG_3_4
%\orbitertableofextracommands{spreads_in_PG_3_4}

\bigskip






Let us illustrate the technique of classification by distinguished substructure.
We consider spreads in $\PG(3,5)$. 
These are sets of lines which partition the set of points. 
The search poset is the set of partial spreads, 
that is, the set of sets of pairwise disjoint lines. 
We decide to pick as substructure the partial spreads consisting of five disjoint lines. 
The number $s=5$ is somewhat arbitrary, chosen only because it seems to work best. 
For each partial spreads of size $s$, the 
lifting constructs all spreads of $\PG(3,5)$ 
containing the chosen partial spread as a subset. 
In a final step, we will perform an isomorph classification on the set of liftings. 
This will furnish the desired classification of spreads of $\PG(3,5).$

\bigskip

Computationally, 
the lifting process is the bottleneck in this procedure. 
To enhance the performance of the lifting, 
we use specialized algorithms from graph theory, based on cliques. 
A clique in a graph is a complete subgraph, 
i.e. a subset of the vertices such that any two 
in the subset are adjacent.
A variant of this is ranbow clique finding. 
Here, we have a coloring of the vertices, 
so that edges are only present between vertices of different color.
The goal of rainbow clique finding 
is to find cliques whose vertices represent all the different colors.  
Many problems in the field of combinatorial designs can be attacked using rainbow cliques.

\bigskip


To illustrate the technique, let us discuss some examples.
Regarding the substructure, 
we choose the parameter $s=5.$ 
This is based on some experimentation.
The command 

\sourcecode{classify_spreads_25_starter_lift_case_0}

classifies the partial spreads of size $s=5$ and 
prepares for the lifting of the first case only. 
In order to prepare for the lifting, a graph is constructed which 
describes the lines that can be added to the first partial spread. 
The vertices of the graph are the lines 
disjoint from the initial set of 5 lines in the partial spread. 
Two vertices are joined by an edge if the associated lines are disjoint. 
The vertices of the graph are colored according to the very first basis vector in the 
generator matrix of the subspace in reduced row echelon form.
In order to find the rainbow cliques in the graph, the command 

\sourcecode{spreads_25_starter_0_cliques}

can be used. 

\bigskip

The command 

\sourcecode{classify_spreads_25_starter_lift_all_cases}

recomputes the partial spreads of size $s=5$ and prepares for the lifting of all 
orbit representatives (there are 28). 
This leads to 28 graphs, each of which is 
written to a file.
The next command performs the rainbow clique finding in each of the 28 graphs:

\sourcecode{spreads_25_starter_cliques}

The resulting cliques are stored in files. 
The command 

\sourcecode{classify_spreads_25_isomorph}

performs the final isomorph rejection on the spreads arising from the 
rainbow cliques in all cases. 
It results in a transversal of the isomorphism classes 
of spreads of $\PG(3,5).$ 
In total, 21 spreads are found. 
This agrees with the results in the literature, 
see~\cite{CzerwinskiOakden92}.






\bigskip




%Table~\ref{tab:spreads:7:2} lists the solid spreads in $\PG(7,2)$ 
%according to their Orbiter catalogue number (OCN).

%\input GRAPHICS/TABLES/spreads_in_PG_7_2
%\orbitertableofextracommands{spreads_in_PG_7_2}

